% =============================================================================
% Chapter 1 --- Overview and corpus layout
% =============================================================================
\chapter{Overview and corpus layout}
\label{ch:overview}

\section{Why this corpus exists}

Agentic AI systems --- LLM-driven agents that plan across multiple steps,
call tools, and act on external systems --- are governed by a rapidly
evolving mix of regulatory frameworks, empirical research, and conformance
tooling. None of these sources alone tells an organisation whether a
specific deployment is defensible. The \corpus{} corpus brings the three
source classes together, adds a machine-checkable link to the code
implementing each requirement, and exposes the result as both structured
JSON and this human-readable PDF.

Three concrete failure modes motivated the corpus:

\begin{itemize}
  \item \textbf{Requirements drift.} Regulatory text is updated faster than
        organisational policies. A machine-readable copy with a hash lets
        governance teams detect a published update and re-run the
        downstream validators.
  \item \textbf{Evidence lives everywhere.} A single requirement such as
        \emph{``bound the agent's tool access''} touches an MCP server, a
        data-product registry, a CI job, a runbook, and a review checklist.
        Without a central mapping, audit evidence walks are painful and
        inconsistent.
  \item \textbf{Conformance tools don't speak the same vocabulary.} AI
        Verify uses plugin IDs, Moonshot uses benchmark IDs, internal
        controls use EUC reference numbers, and product teams use file
        paths. The requirement schema defined here is the translation
        table.
\end{itemize}

\section{Three source classes}

\begin{description}
  \item[Framework --- MAS MGF for Agentic AI.] A 29-page recommended
    framework organised into four pillars: assess and bound risks upfront;
    make humans meaningfully accountable; implement technical controls and
    processes; enable end-user responsibility. It is the highest-normativity
    document in the corpus.
  \item[Index --- 2025 AI Agent Index.] A descriptive inventory of 30
    deployed agentic systems scored on legal, technical, autonomy,
    ecosystem, evaluation, and safety documentation. It supplies
    transparency requirements rather than normative obligations.
  \item[Research --- Ju \& Aral (2026).] A 2{,}234-participant field
    experiment that produced empirical claims about delegation rates, task
    orientation, and diversity collapse. We treat these as observational
    requirements that inform which metrics a deployment should measure.
\end{description}

\section{Two conformance tools}

\begin{description}
  \item[\tool{aiverify} (v2.0).] A plugin-based testing framework. Eleven
    stock plugins implement fairness, robustness, explainability,
    financial-sector (Veritas), reporting, and process-checklist
    algorithms. Plugins are referenced from requirement records via the
    \texttt{verification[].plugin} field.
  \item[\tool{moonshot-cicd} (v1.1.0).] A CI/CD-focused safety evaluation
    tool exposing four benchmark categories --- hallucination, undesirable
    content, data disclosure, adversarial --- behind a single
    \texttt{moonshot run} command. Referenced via
    \texttt{verification[].testId}.
\end{description}

\section{Corpus layout}

\begin{figure}[h]
\centering
\begin{tikzpicture}[
  node distance=6mm and 8mm,
  every node/.style={font=\small},
  box/.style={draw=reggrey,thick,rounded corners=1pt,inner sep=5pt,
              align=center,minimum height=8mm,fill=white},
  dir/.style={box,fill=reggrey!10},
  gen/.style={box,fill=reggreen!10},
  struct/.style={box,fill=regblue!10},
  outp/.style={box,fill=regamber!10},
  arr/.style={-Latex,thick,draw=reggrey!60}
]
\node[dir] (src)      {\textbf{source/}\\ PDFs (untouched)};
\node[gen,right=of src] (text) {\textbf{extracted-text/}\\ page-by-page TXT};
\node[gen,right=of text] (rag) {\textbf{rag/}\\ chunks + embeddings};
\node[struct,below=of text] (struct) {\textbf{structured/}\\ JSON records + schemas \\ + \tool{validate.py}};
\node[outp,below=of struct] (spec) {\textbf{spec/}\\ this PDF (LaTeX)};

\draw[arr] (src) -- (text) node[midway,above,font=\tiny] {pdftotext};
\draw[arr] (text) -- (rag)  node[midway,above,font=\tiny] {chunker};
\draw[arr] (text.south) |- (struct.west);
\draw[arr] (struct) -- (spec);
\end{tikzpicture}
\caption{Pipeline from source PDFs through to the printable specification.}
\label{fig:layout}
\end{figure}

\begin{table}[h]
\centering
\small
\begin{tabularx}{\textwidth}{lX}
\toprule
\textbf{Directory} & \textbf{Purpose} \\
\midrule
\corpus{source/}          & Binary source documents. Never edited; hashed in \corpus{structured/corpus.json}. \\
\corpus{extracted-text/}  & Derived, regeneratable page-by-page plain text. \\
\corpus{rag/}             & RAG chunks + deterministic embedding records (re-upsertable into Qdrant via \corpus{rag/scripts/upsert\_qdrant.py}). \\
\corpus{structured/}      & JSON Schemas, regulation records, requirement records, conformance-tool records, the corpus index, and the validator. \\
\corpus{spec/}            & This human-readable PDF (LaTeX + tectonic). \\
\corpus{aiverify-main/}   & Vendored AI Verify 2.0 monorepo (Apache-2.0). \\
\corpus{moonshot-cicd-main/} & Vendored Moonshot 1.1.0 monorepo (Apache-2.0). \\
\bottomrule
\end{tabularx}
\caption{Top-level directory layout.}
\label{tab:layout}
\end{table}

\section{End-to-end traceability path}

Figure~\ref{fig:traceability} shows the five-step chain that the validator
enforces: from the hashed source PDF, through the extracted text, a
structured regulation record, a normative requirement record with
cross-references, down to the code or test actually implementing it.

\begin{figure}[h]
\centering
\resizebox{\textwidth}{!}{%
\begin{tikzpicture}[
  node distance=4mm and 4mm,
  every node/.style={font=\small},
  stage/.style={draw=reggrey,thick,rounded corners=1pt,inner sep=5pt,
                minimum width=28mm,minimum height=9mm,align=center,fill=white},
  s1/.style={stage,fill=reggrey!10},
  s2/.style={stage,fill=reggreen!10},
  s3/.style={stage,fill=regblue!10},
  s4/.style={stage,fill=regblue!20},
  s5/.style={stage,fill=regamber!15},
  arr/.style={-Latex,thick,draw=reggrey!60}
]
\node[s1] (pdf)   {PDF\\ \tiny hash pinned};
\node[s2,right=of pdf]   (txt)   {Extracted text\\ \tiny re-generatable};
\node[s3,right=of txt]   (reg)   {Regulation record\\ \tiny sections + scope};
\node[s4,right=of reg]   (req)   {Requirement record\\ \tiny impl[] + verify[]};
\node[s5,right=of req]   (code)  {Code / test / tool\\ \tiny repo path};

\draw[arr] (pdf) -- (txt);
\draw[arr] (txt) -- (reg);
\draw[arr] (reg) -- (req);
\draw[arr] (req) -- (code);
\end{tikzpicture}}
\caption{End-to-end traceability chain verified by \corpus{structured/validate.py}.}
\label{fig:traceability}
\end{figure}

\section{What the validator enforces}

\begin{enumerate}
  \item \textbf{Schema conformance.} Every record under the three
        \corpus{structured/} sub-directories
        (\texttt{regulations/}, \texttt{requirements/},
        \texttt{conformance-tools/}) validates against its Draft-07 JSON
        Schema using an offline \texttt{referencing} registry.
  \item \textbf{Corpus integrity.} Every entry in
        \corpus{structured/corpus.json} resolves to a file on disk, and
        every record on disk is listed in the index.
  \item \textbf{Implementation link existence.} For each
        \texttt{implementation[].path}, the validator resolves the path
        against the repository root and fails if it is missing (unless the
        entry is explicitly recorded as \texttt{kind: "none"} marking a
        known gap).
  \item \textbf{Verification link coherence.} Each
        \texttt{verification[].plugin} must be declared in the
        \tool{aiverify} record; each \texttt{verification[].testId} in the
        \tool{moonshot-cicd} record; each \texttt{verification[].path}
        referencing a test must exist on disk.
\end{enumerate}

The validator is the contract: CI, internal audit, and external reviewers
run the same script and see the same results.
\clearpage
